\Chapter{9}

\Verse{1}
{
    \zA{话说秦邦业父子专候贾家人来送上学之信。}
}
{
    \eA{But to return to our story. Mr. Ch'in, the father, and Ch'in Chung, his son,
        only waited until the receipt, by the hands of a servant,}
}
{
}

\Verse{2}
{
    \zA{原来宝玉急于要和秦钟相遇，遂择 了后日一定上学，打发人送了信。}
}
{
    \eA{of a letter from the Chia family about the date on which they were to go to
        school. Indeed, Pao-yü was only too impatient that he and Ch'in Chung should
        come together, and, without loss of time,}
}
{
}

\Verse{3}
{
    \zA{到了这天，宝玉起来时，袭人早已把书笔文物收 拾停妥，坐在床沿上发闷，见宝玉起来，只得伏侍他梳洗。}
}
{
    \eA{he fixed upon two days later as the day upon which they were definitely to
        begin their studies, and he despatched a servant with a letter to this
        effect. On the day appointed, as soon as it was daylight, Pao-yü turned out
        of bed. Hsi Jen had already by that time got books, pencils and all writing
        necessaries in perfect readiness,}
}
{
}

\Verse{4}
{
    \zA{宝玉见他闷闷的，问道： “好姐姐，你怎么又不喜欢了？}
}
{
    \eA{and was sitting on the edge of the bed in a moping mood; but as soon as she
        saw Pao-yü approach, she was constrained to wait upon him in his toilette
        and ablutions.}
}
{
}

\Verse{5}
{
    \zA{难道怕我上学去，撂的你们冷清了不成？”}
}
{
    \eA{Pao-yü, noticing how despondent she was, made it a point to address her. "My
        dear sister," he said, "how is it you aren't again yourself?}
}
{
}

\Verse{6}
{
    \zA{袭人笑 道：“这是那里的话？}
}
{
    \eA{Is it likely that you bear me a grudge for being about to go to school,
        because when I leave you, you'll all feel dull?" Hsi Jen smiled.}
}
{
}

\Verse{7}
{
    \zA{念书是很好的事，不然就潦倒一辈子了，终久怎么样呢？}
}
{
    \eA{"What an ideal" she replied. "Study is a most excellent thing, and without
        it a whole lifetime is a mere waste, and what good comes in the long run?}
}
{
}

\Verse{8}
{
    \zA{但只 一件：只是念书的时候儿想着书，不念的时候儿想着家。}
}
{
    \eA{There's only one thing, which is simply that when engaged in reading your
        books, you should set your mind on your books; and that you should think of
        home when not engaged in reading.}
}
{
}

\Verse{9}
{
    \zA{总别和他们玩闹，碰见老 爷不是玩的。}
}
{
    \eA{Whatever you do, don't romp together with them, for were you to meet our
        master, your father, it will be no joke!}
}
{
}

\Verse{10}
{
    \zA{虽说是奋志要强，那工课宁可少些，一则贪多嚼不烂，二则身子也要 保重。}
}
{
    \eA{Although it's asserted that a scholar must strain every nerve to excel, yet
        it's preferable that the tasks should be somewhat fewer, as, in the first
        place, when one eats too much, one cannot digest it; and,}
}
{
}

\Verse{11}
{
    \zA{这就是我的意思，你好歹体谅些。”}
}
{
    \eA{in the second place, good health must also be carefully attended to. This is
        my view on the subject,}
}
{
}

\Verse{12}
{
    \zA{袭人说一句，宝玉答应一句。}
}
{
    \eA{and you should at all times consider it in practice." While Hsi Jen gave
        utterance to a sentence,}
}
{
}

\Verse{13}
{
    \zA{袭人又道： “大毛儿衣服我也包好了，交给小子们去了。}
}
{
    \eA{Pao-yü nodded his head in sign of approval of that sentence. Hsi Jen then
        went on to speak. "I've also packed up," she continued,}
}
{
}

\Verse{14}
{
    \zA{学里冷，好歹想着添换，比不得家里 有人照顾。}
}
{
    \eA{"your long pelisse, and handed it to the pages to take it over; so mind,
        when it's cold in the school-room, please remember to put on this extra
        clothing,}
}
{
}

\Verse{15}
{
    \zA{脚炉手炉也交出去了，你可逼着他们给你笼上。}
}
{
    \eA{for it's not like home, where you have people to look after you. The
        foot-stove and hand-stove, I've also sent over; and urge that pack of
        lazy-bones to attend to their work,}
}
{
}

\Verse{16}
{
    \zA{那一起懒贼，你不说他 们乐得不动，白冻坏了你。”}
}
{
    \eA{for if you say nothing, they will be so engrossed in their frolics, that
        they'll be loth to move, and let you, all for nothing, take a chill and ruin
        your constitution."}
}
{
}

\Verse{17}
{
    \zA{宝玉道：“你放心，我自己都会调停的。}
}
{
    \eA{"Compose your mind," replied Pao-yü; "when I go out, I know well enough how
        to attend to everything my own self.}
}
{
}

\Verse{18}
{
    \zA{你们也可别 闷死在这屋里，长和林妹妹一处玩玩儿去才好。”}
}
{
    \eA{But you people shouldn't remain in this room, and mope yourselves to death;
        and it would be well if you would often go over to cousin Lin's for a romp."}
}
{
}

\Verse{19}
{
    \zA{说着俱已穿戴齐备，袭人催他去 见贾母、贾政、王夫人。}
}
{
    \eA{While saying this, he had completed his toilette, and Hsi Jen pressed him to
        go and wish good morning to dowager lady Chia, Chia Cheng, madame Wang,}
}
{
}

\Verse{20}
{
    \zA{宝玉又嘱咐了晴雯麝月几句，方出来见贾母。}
}
{
    \eA{and the other members of the family. Pao-yü, after having gone on to give a
        few orders to Ch'ing Wen and She Yueh, at length left his apartments,}
}
{
}

\Verse{21}
{
    \zA{贾母也不免 有几句嘱咐的话。}
}
{
    \eA{and coming over, paid his obeisance to dowager lady Chia. Her venerable
        Ladyship had likewise,}
}
{
}

\Verse{22}
{
    \zA{然后去见王夫人，又出来到书房中见贾政。}
}
{
    \eA{as a matter of course, a few recommendations to make to him, which ended, he
        next went and greeted madame Wang; and leaving again her quarters,}
}
{
}

\Verse{23}
{
    \zA{这日贾政正在书房中和清客相公们说闲话儿，忽见宝玉进来请安，回说上学 去。}
}
{
    \eA{he came into the library to wish Chia Cheng good morning. As it happened,
        Chia Cheng had on this day returned home at an early hour, and was, at this
        moment, in the library, engaged in a friendly chat with a few gentlemen, who
        were family companions.}
}
{
}

\Verse{24}
{
    \zA{贾政冷笑道：“你要再提‘上学’两个字，连我也羞死了。}
}
{
    \eA{Suddenly perceiving Pao-yü come in to pay his respects, and report that he
        was about to go to school, Chia Cheng gave a sardonic smile. "If you do
        again," he remarked,}
}
{
}

\Verse{25}
{
    \zA{依我的话，你竟玩 你的去是正经。}
}
{
    \eA{"make allusions to the words going to school, you'll make even me blush to
        death with shame! My advice to you is that you should after all go your own
        way and play;}
}
{
}

\Verse{26}
{
    \zA{看仔细站腌？}
}
{
    \eA{that's the best thing for you;}
}
{
}

\Verse{27}
{
    \zA{了我这个地，靠腌？}
}
{
    \eA{and mind you don't pollute with dirt this floor by standing here,}
}
{
}

\Verse{28}
{
    \zA{了我这个门！”}
}
{
    \eA{and soil this door of mine by leaning against it!"}
}
{
}

\Verse{29}
{
    \zA{众清客都起身笑 道：“老世翁何必如此。}
}
{
    \eA{The family companions stood up and smilingly expostulated. "Venerable Sir,"
        they pleaded, "why need you be so down upon him?}
}
{
}

\Verse{30}
{
    \zA{今日世兄一去，二三年就可显身成名的，断不似往年仍作 小儿之态了。}
}
{
    \eA{Our worthy brother is this day going to school, and may in two or three
        years be able to display his abilities and establish his reputation. He
        will, beyond doubt, not behave like a child, as he did in years gone past.}
}
{
}

\Verse{31}
{
    \zA{天也将饭时了，世兄竟快请罢。”}
}
{
    \eA{But as the time for breakfast is also drawing nigh, you should, worthy
        brother, go at once." When these words had been spoken,}
}
{
}

\Verse{32}
{
    \zA{说着便有两个年老的携了宝玉出去。}
}
{
    \eA{two among them, who were advanced in years, readily took Pao-yü by the hand,
        and led him out of the library.}
}
{
}

\Verse{33}
{
    \zA{贾政因问：“跟宝玉的是谁？”}
}
{
    \eA{"Who are in attendance upon Pao-yü?" Chia Cheng having inquired, he heard a
        suitable reply,}
}
{
}

\Verse{34}
{
    \zA{只听见外面答应了一声，早进来三四个大汉，打千 儿请安。}
}
{
    \eA{"We, Sir!" given from outside; and three or four sturdy fellows entered at
        an early period and fell on one knee, and bowed and paid their obeisance.
        When Chia Cheng came to scrutinise who they were,}
}
{
}

\Verse{35}
{
    \zA{贾政看时，是宝玉奶姆的儿子名唤李贵的，因向他道：“你们成日家跟他 上学，他到底念了些什么书!倒念了些流言混话在肚子里，学了些精致的淘气。}
}
{
    \eA{and he recognised Li Kuei, the son of Pao-yü's nurse, he addressed himself
        to him. "You people," he said, "remain waiting upon him the whole day long
        at school, but what books has he after all read? Books indeed! why, he has
        read and filled his brains with a lot of trashy words and nonsensical
        phrases,}
}
{
}

\Verse{36}
{
    \zA{等 我闲一闲，先揭了你的皮，再和那不长进的东西算帐！”}
}
{
    \eA{and learnt some ingenious way of waywardness. Wait till I have a little
        leisure, and I'll set to work, first and foremost, and flay your skin off,
        and then settle accounts with that good-for-nothing!"}
}
{
}

\Verse{37}
{
    \zA{吓的李贵忙双膝跪下，摘 了帽子碰头，连连答应“是”，又回说：“哥儿已经念到第三本《诗经》，什么‘攸 攸鹿鸣，荷叶浮萍’，小的不敢撒谎。”}
}
{
    \eA{This threat so terrified Li Kuei that he hastily fell on both his knees,
        pulled off his hat, knocked his head on the ground, and gave vent to
        repeated assenting utterances: "Oh, quite so, Sir! Our elder brother Mr. Pao
        has," he continued, "already read up to the third book of the Book of Odes,}
}
{
}

\Verse{38}
{
    \zA{说的满坐哄然大笑起来，贾政也掌不住笑 了。}
}
{
    \eA{up to where there's something or other like: 'Yiu, Yiu, the deer bleat; the
        lotus leaves and duckweed.' Your servant wouldn't presume to tell a lie!"}
}
{
}

\Verse{39}
{
    \zA{因说道：“那怕再念三十本《诗经》，也是‘掩耳盗铃’，哄人而已。}
}
{
    \eA{As he said this, the whole company burst out into a boisterous fit of
        laughter, and Chia Cheng himself could not also contain his countenance and
        had to laugh. "Were he even," he observed, "to read thirty books of the Book
        of Odes,}
}
{
}

\Verse{40}
{
    \zA{你去请 学里太爷的安，就说我说的：什么《诗经》、古文，一概不用虚应故事，只是先把 《四书》一齐讲明背熟是最要紧的。”}
}
{
    \eA{it would be as much an imposition upon people and no more, as (when the
        thief) who, in order to steal the bell, stops up his own ears! You go and
        present my compliments to the gentleman in the schoolroom, and tell him,
        from my part, that the whole lot of Odes and old writings are of no use, as
        they are subjects for empty show;}
}
{
}

\Verse{41}
{
    \zA{李贵忙答应“是”，见贾政无话，方起来退 出去。}
}
{
    \eA{and that he should, above all things, take the Four Books, and explain them
        to him, from first to last, and make him know them all thoroughly by
        heart,--that this is the most important thing!"}
}
{
}

\Verse{42}
{
    \zA{此时宝玉独站在院外，屏声静候，等他们出来同走。}
}
{
    \eA{Li Kuei signified his obedience with all promptitude, and perceiving that
        Chia Cheng had nothing more to say, he retired out of the room.}
}
{
}

\Verse{43}
{
    \zA{李贵等一面掸衣裳，一面 说道：“哥儿可听见了？}
}
{
    \eA{During this while, Pao-yü had been standing all alone outside in the court,
        waiting quietly with suppressed voice, and when they came out he at once
        walked away in their company.}
}
{
}

\Verse{44}
{
    \zA{先要揭我们的皮呢。}
}
{
    \eA{Li Kuei and his companions observed as they shook their clothes,}
}
{
}

\Verse{45}
{
    \zA{人家的奴才跟主子赚些个体面，我们 这些奴才白陪着挨打受骂的。}
}
{
    \eA{"Did you, worthy brother, hear what he said that he would first of all flay
        our skins off! People's servants acquire some respectability from the master
        whom they serve, but we poor fellows fruitlessly wait upon you,}
}
{
}

\Verse{46}
{
    \zA{从此也可怜见些才好！”}
}
{
    \eA{and are beaten and blown up in the bargain. It would be well if we were,}
}
{
}

\Verse{47}
{
    \zA{宝玉笑道：“好哥哥，你别 委屈，我明儿请你。”}
}
{
    \eA{from henceforward, to be treated with a certain amount of regard." Pao-yü
        smiled, "Dear Brother," he added, "don't feel aggrieved;}
}
{
}

\Verse{48}
{
    \zA{李贵道：“小祖宗，谁敢望‘请’，只求听一两句话就有了。”}
}
{
    \eA{I'll invite you to come round to-morrow!" "My young ancestor," replied Li
        Kuei, "who presumes to look forward to an invitation? all I entreat you is
        to listen to one or two words I have to say,}
}
{
}

\Verse{49}
{
    \zA{说着又至贾母这边，秦钟早已来了，贾母正和他说话儿呢。}
}
{
    \eA{that's all." As they talked they came over once more to dowager lady Chia's
        on this side. Ch'in Chung had already arrived, and the old lady was first
        having a chat with him.}
}
{
}

\Verse{50}
{
    \zA{于是二人见过，辞 了贾母。}
}
{
    \eA{Forthwith the two of them exchanged salutations, and took leave of her
        ladyship;}
}
{
}

\Verse{51}
{
    \zA{宝玉忽想起未辞黛玉，又忙至黛玉房中来作辞。}
}
{
    \eA{but Pao-yü, suddenly remembering that he had not said good-bye to Tai-yü,
        promptly betook himself again to Tai-yü's quarters to do so. Tai-yü was, at
        this time,}
}
{
}

\Verse{52}
{
    \zA{彼时黛玉在窗下对镜理妆， 听宝玉说上学去，因笑道：“好!这一去，可是要‘蟾宫折桂’了!我不能送你了。”}
}
{
    \eA{below the window, facing the mirror, and adjusting her toilette. Upon
        hearing Pao-yü mention that he was on his way to school, she smiled and
        remarked, "That's right! you're now going to school and you'll be sure to
        reach the lunar palace and pluck the olea fragrans;}
}
{
}

\Verse{53}
{
    \zA{宝玉道：“好妹妹，等我下学再吃晚饭。}
}
{
    \eA{but I can't go along with you." "My dear cousin," rejoined Pao-yü, "wait for
        me to come out from school,}
}
{
}

\Verse{54}
{
    \zA{那胭脂膏子也等我来再制。”}
}
{
    \eA{before you have your evening meal; wait also until I come to prepare the
        cosmetic of rouge."}
}
{
}

\Verse{55}
{
    \zA{唠叨了半日， 方抽身去了。}
}
{
    \eA{After a protracted chat, he at length tore himself away and took his
        departure.}
}
{
}

\Verse{56}
{
    \zA{黛玉忙又叫住，问道：“你怎么不去辞你宝姐姐来呢？”}
}
{
    \eA{"How is it," interposed Tai-yü, as she once again called out to him and
        stopped him, "that you don't go and bid farewell to your cousin Pao Ch'ai?"}
}
{
}

\Verse{57}
{
    \zA{宝玉笑而不 答，一径同秦钟上学去了。}
}
{
    \eA{Pao-yü smiled, and saying not a word by way of reply he straightway walked
        to school, accompanied by Ch'in Chung.}
}
{
}

\Verse{58}
{
    \zA{原来这义学也离家不远，原系当日始祖所立，恐族中子弟有力不能延师者，即 入此中读书。}
}
{
    \eA{This public school, which it must be noticed was also not far from his
        quarters, had been originally instituted by the founder of the
        establishment, with the idea that should there be among the young fellows of
        his clan any who had not the means to engage a tutor,}
}
{
}

\Verse{59}
{
    \zA{凡族中为官者皆有帮助银两，以为学中膏火之费；举年高有德之人为 塾师。}
}
{
    \eA{they should readily be able to enter this class for the prosecution of their
        studies; that all those of the family who held official position should all
        give (the institution) pecuniary assistance, with a view to meet the
        expenses necessary for allowances to the students;}
}
{
}

\Verse{60}
{
    \zA{如今秦宝二人来了，一一的都互相拜见过，读起书来。}
}
{
    \eA{and that they were to select men advanced in years and possessed of virtue
        to act as tutors of the family school. The two of them, Ch'in Chung and
        Pao-yü, had now entered the class,}
}
{
}

\Verse{61}
{
    \zA{自此后二人同来同往 同起同坐，愈加亲密。}
}
{
    \eA{and after they and the whole number of their schoolmates had made each
        other's acquaintance, their studies were commenced.}
}
{
}

\Verse{62}
{
    \zA{兼贾母爱惜，也常留下秦钟一住三五天，和自己重孙一般看 待。}
}
{
    \eA{Ever since this time, these two were wont to come together, go together, get
        up together, and sit together, till they became more intimate and close.
        Besides, dowager lady Chia got very fond of Ch'in Chung,}
}
{
}

\Verse{63}
{
    \zA{因见秦钟家中不甚宽裕，又助些衣服等物。}
}
{
    \eA{and would again and again keep him to stay with them for three and five days
        at a time, treating him as if he were one of her own great-grandsons.}
}
{
}

\Verse{64}
{
    \zA{不上一两月工夫，秦钟在荣府里便 惯熟了。}
}
{
    \eA{Perceiving that in Ch'in Chung's home there was not much in the way of
        sufficiency, she also helped him in clothes and other necessaries;}
}
{
}

\Verse{65}
{
    \zA{宝玉终是个不能安分守理的人，一味的随心所欲，因此发了癖性，又向秦 钟悄说：“咱们两个人，一样的年纪，况又同窗，以后不必论叔侄，只论弟兄朋友 就是了。”}
}
{
    \eA{and scarcely had one or two months elapsed before Ch'in Chung got on
        friendly terms with every one in the Jung mansion. Pao-yü was, however, a
        human being who could not practise contentment and observe propriety; and as
        his sole delight was to have every caprice gratified, he naturally developed
        a craving disposition. "We two, you and I, are," he was also wont secretly
        to tell Ch'in Chung, "of the same age, and fellow-scholars besides,}
}
{
}

\Verse{66}
{
    \zA{先是秦钟不敢，宝玉不从，只叫他“兄弟”，叫他表字“鲸卿”，秦钟 也只得混着乱叫起来。}
}
{
    \eA{so that there's no need in the future to pay any regard to our relationship
        of uncle and nephew; and we should treat each other as brothers or friends,
        that's all." Ch'in Chung at first (explained that) he could not be so
        presumptuous; but as Pao-yü would not listen to any such thing,}
}
{
}

\Verse{67}
{
    \zA{原来这学中虽都是本族子弟与些亲戚家的子侄，俗语说的好：“一龙九种，种 种各别。”}
}
{
    \eA{but went on to address him as brother and to call him by his style Ch'ing
        Ch'ing, he had likewise himself no help, but to begin calling him, at
        random, anything and anyhow. There were, it is true, a large number of
        pupils in this school,}
}
{
}

\Verse{68}
{
    \zA{未免人多了就有龙蛇混杂、下流人物在内。}
}
{
    \eA{but these consisted of the sons and younger brothers of that same clan, and
        of several sons and nephews of family connections.}
}
{
}

\Verse{69}
{
    \zA{自秦宝二人来了，都生的花 朵儿一般的模样，又见秦钟腼腆温柔，未语先红，怯怯羞羞有女儿之风；宝玉又是 天生成惯能作小服低，赔身下气，性情体贴，话语缠绵。}
}
{
    \eA{The proverb appositely describes that there are nine species of dragons, and
        that each species differs; and it goes of course without saying that in a
        large number of human beings there were dragons and snakes, confusedly
        admixed, and that creatures of a low standing were included. Ever since the
        arrival of the two young fellows, Ch'in Chung and Pao-yü, both of whom were
        in appearance as handsome as budding flowers, and they, on the one hand,}
}
{
}

\Verse{70}
{
    \zA{因他二人又这般亲厚，也 怨不得那起同窗人起了嫌疑之念，背地里你言我语，诟谇谣诼，布满书房内外。}
}
{
    \eA{saw how modest and genial Ch'in Chung was, how he blushed before he uttered
        a word, how he was timid and demure like a girl, and on the other hand, how
        that Pao-yü was naturally proficient in abasing and demeaning himself, how
        he was so affable and good-natured, considerate in his temperament and so
        full of conversation,}
}
{
}

\Verse{71}
{
    \zA{原来薛蟠自来王夫人处住后，便知有一家学，学中广有青年子弟。}
}
{
    \eA{and how that these two were, in consequence, on such terms of intimate
        friendship, it was, in fact, no matter of surprise that the whole company of
        fellow-students began to foster envious thoughts,}
}
{
}

\Verse{72}
{
    \zA{偶动了龙阳 之兴，因此也假说来上学，不过是“三日打鱼，两日晒网”，白送些束？}
}
{
    \eA{that they, behind their backs, passed on their account, this one one
        disparaging remark and that one another, and that they insinuated slanderous
        lies against them, which extended inside as well as outside the school-room.}
}
{
}

\Verse{73}
{
    \zA{礼物与贾 代儒，却不曾有一点儿进益，只图结交些契弟。}
}
{
    \eA{Indeed, after Hsüeh P'an had come over to take up his quarters in madame
        Wang's suite of apartments, he shortly came to hear of the existence of a
        family school, and that this school was mainly attended by young fellows of
        tender years,}
}
{
}

\Verse{74}
{
    \zA{谁想这学内的小学生，图了薛蟠的 银钱穿吃，被他哄上手了，也不消多记。}
}
{
    \eA{and inordinate ideas were suddenly aroused in him. While he therefore
        fictitiously gave out that he went to school, [he was as irregular in his
        attendance as the fisherman] who catches fish for three days, and suns his
        nets for the next two;}
}
{
}

\Verse{75}
{
    \zA{又有两个多情的小学生，亦不知是那一房 的亲眷，亦未考真姓名，只因生得妩媚风流，满学中都送了两个外号，一个叫“香 怜”，一个叫“玉爱”。}
}
{
    \eA{simply presenting his school-fee gift to Chia Tai-jui and making not the
        least progress in his studies; his sole dream being to knit a number of
        familiar friendships. Who would have thought it, there were in this school
        young pupils, who, in their greed to obtain money, clothes and eatables from
        Hsüeh P'an, allowed themselves to be cajoled by him,}
}
{
}

\Verse{76}
{
    \zA{别人虽都有羡慕之意、“不利于孺子”之心，只是惧怕薛 蟠的威势，不敢来沾惹。}
}
{
    \eA{and played tricks upon; but on this topic, it is likewise superfluous to
        dilate at any length. There were also two lovable young scholars, relatives
        of what branch of the family is not known, and whose real surnames and names
        have also not been ascertained,}
}
{
}

\Verse{77}
{
    \zA{如今秦宝二人一来了，见了他两个，也不免缱绻羡爱，亦 知系薛蟠相知，未敢轻举妄动。}
}
{
    \eA{who, by reason of their good and winsome looks, were, by the pupils in the
        whole class, given two nicknames, to one that of "Hsiang Lin," "Fragrant
        Love," and to the other "Yü Ai," "Precious Affection." But although every
        one entertained feelings of secret admiration for them,}
}
{
}

\Verse{78}
{
    \zA{香玉二人心中，一般的留情与秦宝：因此四人心中 虽有情意，只未发出。}
}
{
    \eA{and had the wish to take liberties with the young fellows, they lived,
        nevertheless, one and all, in such terror of Hsüeh P'an's imperious
        influence, that they had not the courage to come forward and interfere with
        them.}
}
{
}

\Verse{79}
{
    \zA{每日一入学中，四处各坐，却八目勾留，或设言托意，或咏 桑寓柳，遥以心照，却外面自为避人眼目。}
}
{
    \eA{As soon as Ch'in Chung and Pao-yü had, at this time, come to school, and
        they had made the acquaintance of these two fellow-pupils, they too could
        not help becoming attached to them and admiring them, but as they also came
        to know that they were great friends of Hsüeh P'an,}
}
{
}

\Verse{80}
{
    \zA{不料偏又有几个滑贼看出形景来，都背 后挤眉弄眼，或咳嗽扬声，这也非止一日。}
}
{
    \eA{they did not, in consequence, venture to treat them lightly, or to be
        unseemly in their behaviour towards them. Hsiang Lin and Yü Ai both kept to
        themselves the same feelings, which they fostered for Ch'in Chung and
        Pao-yü,}
}
{
}

\Verse{81}
{
    \zA{可巧这日代儒有事回家，只留下一句七言对联，令学生对了明日再来上书，将 学中之事又命长孙贾瑞管理。}
}
{
    \eA{and to this reason is to be assigned the fact that though these four persons
        nurtured fond thoughts in their hearts there was however no visible sign of
        them. Day after day, each one of them would, during school hours, sit in
        four distinct places: but their eight eyes were secretly linked together;
        and, while indulging either in innuendoes or in double entendres,}
}
{
}

\Verse{82}
{
    \zA{妙在薛蟠如今不大上学应卯了，因此秦钟趁此和香怜 弄眉挤眼，二人假出小恭，走至后院说话。}
}
{
    \eA{their hearts, in spite of the distance between them, reflected the whole
        number of their thoughts. But though their outward attempts were devoted to
        evade the detection of other people's eyes, it happened again that, while
        least expected, several sly lads discovered the real state of affairs,}
}
{
}

\Verse{83}
{
    \zA{秦钟先问他：“家里的大人可管你交朋 友不管？”}
}
{
    \eA{with the result that the whole school stealthily frowned their eyebrows at
        them, winked their eyes at them, or coughed at them, or raised their voices
        at them;}
}
{
}

\Verse{84}
{
    \zA{一语未了，只听见背后咳嗽了一声。}
}
{
    \eA{and these proceedings were, in fact, not restricted to one single day. As
        luck would have it, on this day Tai-jui was,}
}
{
}

\Verse{85}
{
    \zA{二人吓的忙回顾时，原来是窗友名 金荣的。}
}
{
    \eA{on account of business, compelled to go home; and having left them as a task
        no more than a heptameter line for an antithetical couplet, explaining that
        they should find a sentence to rhyme,}
}
{
}

\Verse{86}
{
    \zA{香怜本有些性急，便羞怒相激，问他道：“你咳嗽什么？}
}
{
    \eA{and that the following day when he came back, he would set them their
        lessons, he went on to hand the affairs connected with the class to his
        elder grandson, Chia Jui, whom he asked to take charge.}
}
{
}

\Verse{87}
{
    \zA{难道不许我们说 话不成？”}
}
{
    \eA{Wonderful to say Hsüeh P'an had of late not frequented school very often,
        not even so much as to answer the roll,}
}
{
}

\Verse{88}
{
    \zA{金荣笑道：“许你们说话，难道不许我咳嗽不成？}
}
{
    \eA{so that Ch'in Chung availed himself of his absence to ogle and smirk with
        Hsiang Lin; and these two pretending that they had to go out, came into the
        back court for a chat.}
}
{
}

\Verse{89}
{
    \zA{我只问你们：有话不 分明说，许你们这样鬼鬼祟祟的干什么故事？}
}
{
    \eA{"Does your worthy father at home mind your having any friends?" Ch'in Chung
        was the first to ask. But this sentence was scarcely ended, when they heard
        a sound of coughing coming from behind.}
}
{
}

\Verse{90}
{
    \zA{我可也拿住了!还赖什么？}
}
{
    \eA{Both were taken much aback, and, speedily turning their heads round to see,}
}
{
}

\Verse{91}
{
    \zA{先让我抽个 头儿，咱们一声儿不言语。}
}
{
    \eA{they found that it was a fellow-scholar of theirs, called Chin Jung. Hsiang
        Lin was naturally of somewhat hasty temperament,}
}
{
}

\Verse{92}
{
    \zA{不然大家就翻起来！”}
}
{
    \eA{so that with shame and anger mutually impelling each other,}
}
{
}

\Verse{93}
{
    \zA{秦香二人就急得飞红的脸，便问 道：“你拿住什么了？”}
}
{
    \eA{he inquired of him, "What's there to cough at? Is it likely you wouldn't
        have us speak to each other?" "I don't mind your speaking," Chin Jung
        observed laughing;}
}
{
}

\Verse{94}
{
    \zA{金荣笑道：“我现拿住了是真的。”}
}
{
    \eA{"but would you perchance not have me cough? I'll tell you what, however; if
        you have anything to say,}
}
{
}

\Verse{95}
{
    \zA{说着又拍着手笑嚷道： “贴的好烧饼!你们都不买一个吃去？”}
}
{
    \eA{why not utter it in intelligible language? Were you allowed to go on in this
        mysterious manner, what strange doings would you be up to? But I have sure
        enough found you out,}
}
{
}

\Verse{96}
{
    \zA{秦钟香怜二人又气又急，忙进来向贾瑞前 告金荣，说金荣无故欺负他两个。}
}
{
    \eA{so what's the need of still prevaricating? But if you will, first of all,
        let me partake of a share in your little game, you and I can hold our tongue
        and utter not a word. If not, why the whole school will begin to turn the
        matter over."}
}
{
}

\Verse{97}
{
    \zA{原来这贾瑞最是个图便宜没行止的人，每在学中以公报私，勒索子弟们请他； 后又助着薛蟠图些银钱酒肉，一任薛蟠横行霸道，他不但不去管约，反助纣为虐讨 好儿。}
}
{
    \eA{At these words, Ch'in Chung and Hsiang Lin were so exasperated that their
        blood rushed up to their faces. "What have you found out?" they hastily
        asked. "What I have now detected," replied Chin Jung smiling, "is the plain
        truth!" and saying this he went on to clap his hands and to call out with a
        loud voice as he laughed: "They have moulded some nice well-baked cakes,
        won't you fellows come and buy one to eat!"}
}
{
}

\Verse{98}
{
    \zA{偏那薛蟠本是浮萍心性，今日爱东，明日爱西，近来有了新朋友，把香玉二 人丢开一边；就连金荣也是当日的好友，自有了香玉二人，便见弃了金荣，近日连
        香玉亦已见弃。}
}
{
    \eA{(These two have been up to larks, won't you come and have some fun!) Both
        Ch'in Chung and Hsiang Lin felt resentful as well as fuming with rage, and
        with hurried step they went in, in search of Chia Jui, to whom they reported
        Chin Jung, explaining that Chin Jung had insulted them both, without any
        rhyme or reason.}
}
{
}

\Verse{99}
{
    \zA{故贾瑞也无了提携帮衬之人，不怨薛蟠得新厌故，只怨香玉二人不 在薛蟠跟前提携了：因此贾瑞金荣等一干人，也正醋妒他两个。}
}
{
    \eA{The fact is that this Chia Jui was, in an extraordinary degree, a man with
        an eye to the main chance, and devoid of any sense of propriety. His wont
        was at school to take advantage of public matters to serve his private
        interest, and to bring pressure upon his pupils with the intent that they
        should regale him. While subsequently he also lent his countenance to Hsüeh
        P'an,}
}
{
}

\Verse{100}
{
    \zA{今见秦香二人来告 金荣，贾瑞心中便不自在起来，虽不敢呵叱秦钟，却拿着香怜作法，反说他多事， 着实抢白了几句。}
}
{
    \eA{scheming to get some money or eatables out of him, he left him entirely free
        to indulge in disorderly behaviour; and not only did he not go out of his
        way to hold him in check, but, on the contrary, he encouraged him, infamous
        though he was already, to become a bully, so as to curry favour with him.
        But this Hsüeh P'an was,}
}
{
}

\Verse{101}
{
    \zA{香怜反讨了没趣，连秦钟也讪讪的各归坐位去了。}
}
{
    \eA{by nature, gifted with a fickle disposition; to-day, he would incline to the
        east, and to-morrow to the west, so that having recently obtained new
        friends,}
}
{
}

\Verse{102}
{
    \zA{金荣越发得了意，摇头咂嘴的，口内还说许多闲话。}
}
{
    \eA{he put Hsiang Lin and Yü Ai aside. Chin Jung too was at one time an intimate
        friend of his, but ever since he had acquired the friendship of the two
        lads, Hsiang Lin and Yü Ai, he forthwith deposed Chin Jung. Of late, he had
        already come to look down upon even Hsiang Lin and Yü Ai,}
}
{
}

\Verse{103}
{
    \zA{玉爱偏又听见，两个人隔 坐咕咕唧唧的角起口来。}
}
{
    \eA{with the result that Chia Jui as well was deprived of those who could lend
        him support, or stand by him; but he bore Hsüeh P'an no grudge, for wearying
        with old friends, as soon as he found new ones,}
}
{
}

\Verse{104}
{
    \zA{金荣只一口咬定说：“方才明明的撞见他两个在后院里亲 嘴摸屁股，两个商议，定了一对儿。”}
}
{
    \eA{but felt angry that Hsiang Lin and Yü Ai had not put in a word on his behalf
        with Hsüeh P'an. Chia Jui, Chin Jung and in fact the whole crowd of them
        were, for this reason, just harbouring a jealous grudge against these two,
        so that when he saw Ch'in Chung and Hsiang Lin come on this occasion and
        lodge a complaint against Chin Jung,}
}
{
}

\Verse{105}
{
    \zA{论长道短，那时只顾得志乱说，却不防还有 别人。}
}
{
    \eA{Chia Jui readily felt displeasure creep into his heart; and, although he did
        not venture to call Ch'in Chung to account, he nevertheless made an example
        of Hsiang Lin.}
}
{
}

\Verse{106}
{
    \zA{谁知早又触怒了一个人。}
}
{
    \eA{And instead (of taking his part), he called him a busybody and denounced him
        in much abusive language,}
}
{
}

\Verse{107}
{
    \zA{你道这一个人是谁？}
}
{
    \eA{with the result that Hsiang Lin did not, contrariwise, profit in any way,}
}
{
}

\Verse{108}
{
    \zA{原来这人名唤贾蔷，亦系宁府 中之正派玄孙，父母早亡，从小儿跟着贾珍过活，如今长了十六岁，比贾蓉生得还 风流俊俏。}
}
{
    \eA{but brought displeasure upon himself. Even Ch'in Chung grumbled against the
        treatment, as each of them resumed their places. Chin Jung became still more
        haughty, and wagging his head and smacking his lips, he gave vent to many
        more abusive epithets; but as it happened that they also reached Yü Ai's
        ears,}
}
{
}

\Verse{109}
{
    \zA{他兄弟二人最相亲厚，常共起居，宁府中人多口杂，那些不得志的奴仆， 专能造言诽谤主人，因此不知又有什么小人诟谇谣诼之辞。}
}
{
    \eA{the two of them, though seated apart, began an altercation in a loud tone of
        voice. Chin Jung, with obstinate pertinacity, clung to his version. "Just a
        short while back," he said, "I actually came upon them, as they were
        indulging in demonstrations of intimate friendship in the back court. These
        two had resolved to be one in close friendship,}
}
{
}

\Verse{110}
{
    \zA{贾珍想亦风闻得些口声 不好，自己也要避些嫌疑，如今竟分与房舍，命贾蔷搬出宁府，自己立门户过活去 了。}
}
{
    \eA{and were eloquent in their protestations, mindful only in persistently
        talking their trash, but they were not aware of the presence of another
        person." But his language had, contrary to all expectations, given, from the
        very first, umbrage to another person, and who do you, (gentle reader,)}
}
{
}

\Verse{111}
{
    \zA{这贾蔷外相既美，内性又聪敏，虽然应名来上学，亦不过虚掩眼目而已，仍是 斗鸡走狗、赏花阅柳为事。}
}
{
    \eA{imagine this person to have been? This person was, in fact, one whose name
        was Chia Se; a grandson likewise of a main branch of the Ning mansion. His
        parents had died at an early period, and he had, ever since his youth, lived
        with Chia Chen. He had at this time grown to be sixteen years of age,}
}
{
}

\Verse{112}
{
    \zA{上有贾珍溺爱，下有贾蓉匡助，因此族中人谁敢触逆于 他。}
}
{
    \eA{and was, as compared with Chia Jung, still more handsome and good looking.
        These two cousins were united by ties of the closest intimacy, and were
        always together, whether they went out or stayed at home.}
}
{
}

\Verse{113}
{
    \zA{他既和贾蓉最好，今见有人欺负秦钟，如何肯依？}
}
{
    \eA{The inmates of the Ning mansion were many in number, and their opinions of a
        mixed kind; and that whole bevy of servants, devoid as they were of all
        sense of right,}
}
{
}

\Verse{114}
{
    \zA{如今自己要挺身出来报不平， 心中且忖度一番：“金荣贾瑞一等人，都是薛大叔的相知，我又与薛大叔相好，倘 或我一出头，他们告诉了老薛，我们岂不伤和气呢。}
}
{
    \eA{solely excelled in the practice of inventing stories to backbite their
        masters; and this is how some mean person or other again, who it was is not
        known, insinuated slanderous and opprobrious reports (against Chia Se). Chia
        Chen had, presumably, also come to hear some unfavourable criticisms (on his
        account), and having, of course, to save himself from odium and suspicion,}
}
{
}

\Verse{115}
{
    \zA{欲要不管，这谣言说的大家没 趣。}
}
{
    \eA{he had, at this juncture, after all, to apportion him separate quarters, and
        to bid Chia Se move outside the Ning mansion,}
}
{
}

\Verse{116}
{
    \zA{如今何不用计制伏，又止息了口声，又不伤脸面。”}
}
{
    \eA{where he went and established a home of his own to live in. This Chia Se was
        handsome as far as external appearances went, and intelligent withal in his
        inward natural gifts, but, though he nominally came to school, it was simply
        however as a mere blind;}
}
{
}

\Verse{117}
{
    \zA{想毕，也装出小恭去，走 至后面瞧瞧，把跟宝玉书童茗烟叫至身边，如此这般，调拨他几句。}
}
{
    \eA{for he treated, as he had ever done, as legitimate occupations, such things
        as cock fighting, dog-racing and visiting places of easy virtue. And as,
        above, he had Chia Chen to spoil him by over-indulgence; and below, there
        was Chia Jung to stand by him,}
}
{
}

\Verse{118}
{
    \zA{这茗烟乃是宝玉第一个得用且又年轻不谙事的，今听贾蔷说：“金荣如此欺负 秦钟，连你们的爷宝玉都干连在内，不给他个知道，下次越发狂纵。”}
}
{
    \eA{who of the clan could consequently presume to run counter to him? Seeing
        that he was on the closest terms of friendship with Chia Jung, how could he
        reconcile himself to the harsh treatment which he now saw Ch'in Chung
        receive from some persons? Being now bent upon pushing himself forward to
        revenge the injustice, he was, for the time, giving himself up to communing
        with his own heart. "Chin Jung, Chia Jui and the rest are,"}
}
{
}

\Verse{119}
{
    \zA{这茗烟无故 就要欺压人的，如今得了这信，又有贾蔷助着，便一头进来找金荣。}
}
{
    \eA{he pondered, "friends of uncle Hsüeh, but I too am on friendly terms with
        him, and he with me, and if I do come forward and they tell old Hsüeh, won't
        we impair the harmony which exists between us? and if I don't concern
        myself,}
}
{
}

\Verse{120}
{
    \zA{也不叫“金相 公”了，只说：“姓金的，你什么东西！”}
}
{
    \eA{such idle tales make, when spoken, every one feel uncomfortable; and why
        shouldn't I now devise some means to hold them in check, so as to stop their
        mouths,}
}
{
}

\Verse{121}
{
    \zA{贾蔷遂跺一跺靴子，故意整整衣服，看 看日影儿说：“正时候了。”}
}
{
    \eA{and prevent any loss of face!" Having concluded this train of thought, he
        also pretended that he had to go out, and, walking as far as the back, he,
        with low voice, called to his side Ming Yen,}
}
{
}

\Verse{122}
{
    \zA{遂先向贾瑞说有事要早走一步。}
}
{
    \eA{the page attending upon Pao-yü in his studies, and in one way and another,
        he made use of several remarks to egg him on.}
}
{
}

\Verse{123}
{
    \zA{贾瑞不敢止他，只得 随他去了。}
}
{
    \eA{This Ming Yen was the smartest of Pao-yü's attendants, but he was also young
        in years and lacked experience,}
}
{
}

\Verse{124}
{
    \zA{这里茗烟走进来，便一把揪住金荣问道：“我们？}
}
{
    \eA{so that he lent a patient ear to what Chia Se had to say about the way Chin
        Jung had insulted Ch'in Chung. "Even your own master, Pao-yü," (Chia Se
        added), "is involved,}
}
{
}

\Verse{125}
{
    \zA{屁股不？}
}
{
    \eA{and if you don't let him know a bit of your mind,}
}
{
}

\Verse{126}
{
    \zA{，管你？？}
}
{
    \eA{he will next time be still more arrogant."}
}
{
}

\Verse{127}
{
    \zA{相干？}
}
{
    \eA{This Ming Yen was always ready,}
}
{
}

\Verse{128}
{
    \zA{横竖没？}
}
{
    \eA{even with no valid excuse,}
}
{
}

\Verse{129}
{
    \zA{你的爹罢了!说你是好小子，出来动一动你茗大爷！”}
}
{
    \eA{to be insolent and overbearing to people, so that after hearing the news and
        being furthermore instigated by Chia Se, he speedily rushed into the
        schoolroom and cried out "Chin Jung;"}
}
{
}

\Verse{130}
{
    \zA{吓的满屋中子弟都 忙忙的痴望。}
}
{
    \eA{nor did he address him as Mr. Chin, but merely shouted "What kind of fellow
        is this called Chin?"}
}
{
}

\Verse{131}
{
    \zA{贾瑞忙喝：“茗烟不得撒野！”}
}
{
    \eA{Chia Se presently shuffled his feet, while he designedly adjusted his dress
        and looked at the rays of the sun.}
}
{
}

\Verse{132}
{
    \zA{金荣气黄了脸，说：“反了!奴才小 子都敢如此，我只和你主子说。”}
}
{
    \eA{"It's time," he observed and walking forthwith, first up to Chia Jui, he
        explained to him that he had something to attend to and would like to get
        away a little early; and as Chia Jui did not venture to stop him, he had no
        alternative but to let him have his way and go. During this while, Ming Yen
        had entered the room and promptly seizing Chin Jung in a grip:}
}
{
}

\Verse{133}
{
    \zA{便夺手要去抓打宝玉。}
}
{
    \eA{"What we do, whether proper or improper," he said, "doesn't concern you!}
}
{
}

\Verse{134}
{
    \zA{秦钟刚转出身来，听得脑 后飕的一声，早见一方砚瓦飞来，并不知系何人打来，却打了贾蓝贾菌的座上。}
}
{
    \eA{It's enough anyway that we don't defile your father! A fine brat you are
        indeed, to come out and meddle with your Mr. Ming!" These words plunged the
        scholars of the whole class in such consternation that they all wistfully
        and absently looked at him. "Ming Yen," hastily shouted out Chia Jui,}
}
{
}

\Verse{135}
{
    \zA{这 贾蓝贾菌亦系荣府近派的重孙。}
}
{
    \eA{"you're not to kick up a rumpus." Chin Jung was so full of anger that his
        face was quite yellow. "What a subversion of propriety!}
}
{
}

\Verse{136}
{
    \zA{这贾菌少孤，其母疼爱非常，书房中与贾蓝最好， 所以二人同坐。}
}
{
    \eA{a slave and a menial to venture to behave in this manner! I'll just simply
        speak to your master," he exclaimed as he readily pushed his hands off and
        was about to go and lay hold of Pao-yü to beat him.}
}
{
}

\Verse{137}
{
    \zA{谁知这贾菌年纪虽小，志气最大，极是淘气不怕人的。}
}
{
    \eA{Ch'in Chung was on the point of turning round to leave the room, when with a
        sound of 'whiff' which reached him from behind, he at once caught sight of a
        square inkslab come flying that way.}
}
{
}

\Verse{138}
{
    \zA{他在位上， 冷眼看见金荣的朋友暗助金荣，飞砚来打茗烟，偏打错了落在自己面前，将个磁砚 水壶儿打粉碎，溅了一书墨水。}
}
{
    \eA{Who had thrown it he could not say, but it struck the desk where Chia Lan
        and Chia Chün were seated. These two, Chia Lan and Chia Chün, were also the
        great-grandsons of a close branch of the Jung mansion. This Chia Chün had
        been left fatherless at an early age,}
}
{
}

\Verse{139}
{
    \zA{贾菌如何依得，便骂：“好囚攮的们!这不都动了 手了么！”}
}
{
    \eA{and his mother doated upon him in an unusual manner, and it was because at
        school he was on most friendly terms with Chia Lan, that these two sat
        together at the same desk.}
}
{
}

\Verse{140}
{
    \zA{骂着，也便抓起砚台来要飞。}
}
{
    \eA{Who would have believed that Chia Chün would, in spite of being young in
        years, have had an extremely strong mind,}
}
{
}

\Verse{141}
{
    \zA{贾蓝是个省事的，忙按住砚台，忙劝道： “好兄弟，不与咱们相干。”}
}
{
    \eA{and that he would be mostly up to mischief without the least fear of any
        one. He watched with listless eye from his seat Chin Jung's friends
        stealthily assist Chin Jung, as they flung an inkslab to strike Ming Yen,}
}
{
}

\Verse{142}
{
    \zA{贾菌如何忍得住，见按住砚台，他便两手抱起书箧子 来照这边扔去。}
}
{
    \eA{but when, as luck would have it, it hit the wrong mark, and fell just in
        front of him, smashing to atoms the porcelain inkslab and water bottle, and
        smudging his whole book with ink, Chia Chün was, of course,}
}
{
}

\Verse{143}
{
    \zA{终是身小力薄，却扔不到，反扔到宝玉秦钟案上就落下来了。}
}
{
    \eA{much incensed, and hastily gave way to abuse. "You consummate pugnacious
        criminal rowdies! why, doesn't this amount to all of you taking a share in
        the fight!" And as he uttered this abuse,}
}
{
}

\Verse{144}
{
    \zA{只听 豁啷一响，砸在桌上，书本、纸片、笔、砚等物撒了一桌，又把宝玉的一碗茶也砸 得碗碎茶流。}
}
{
    \eA{he too forthwith seized an inkslab, which he was bent upon flinging. Chia
        Lan was one who always tried to avoid trouble, so that he lost no time in
        pressing down the inkslab, while with all the words his mouth could express,
        he tried to pacify him, adding "My dear brother,}
}
{
}

\Verse{145}
{
    \zA{那贾菌即便跳出来，要揪打那飞砚的人。}
}
{
    \eA{it's no business of yours and mine." Chia Chün could not repress his
        resentment; and perceiving that the inkslab was held down,}
}
{
}

\Verse{146}
{
    \zA{金荣此时随手抓了一根毛竹 大板在手，地狭人多，那里经得舞动长板。}
}
{
    \eA{he at once laid hold of a box containing books, which he flung in this
        direction; but being, after all, short of stature, and weak of strength, he
        was unable to send it anywhere near the mark; so that it dropped instead
        when it got as far as the desk belonging to Pao-yü and Ch'in Chung,}
}
{
}

\Verse{147}
{
    \zA{茗烟早吃了一下，乱嚷：“你们还不来 动手？”}
}
{
    \eA{while a dreadful crash became audible as it fell smash on the table. The
        books, papers, pencils, inkslabs, and other writing materials were all
        scattered over the whole table;}
}
{
}

\Verse{148}
{
    \zA{宝玉还有几个小厮，一名扫红，一名锄药，一名墨雨，这三个岂有不淘气 的，一齐乱嚷：“小妇养的!动了兵器了！”}
}
{
    \eA{and Pao-yü's cup besides containing tea was itself broken to pieces and the
        tea spilt. Chia Chün forthwith jumped forward with the intent of assailing
        the person who had flung the inkslab at the very moment that Chin Jung took
        hold of a long bamboo pole which was near by; but as the space was limited,
        and the pupils many, how could he very well brandish a long stick? Ming Yen
        at an early period received a whack,}
}
{
}

\Verse{149}
{
    \zA{墨雨遂掇起一根门闩，扫红锄药手中 都是马鞭子，蜂拥而上。}
}
{
    \eA{and he shouted wildly, "Don't you fellows yet come to start a fight." Pao-yü
        had, besides, along with him several pages, one of whom was called Sao Hung,
        another Ch'u Yo, another Mo Yü.}
}
{
}

\Verse{150}
{
    \zA{贾瑞急得拦一回这个，劝一回那个，谁听他的话？}
}
{
    \eA{These three were naturally up to every mischief, so that with one voice,
        bawling boisterously, "You children of doubtful mothers, have you taken up
        arms?"}
}
{
}

\Verse{151}
{
    \zA{肆行大 乱。}
}
{
    \eA{Mo Yü promptly took up the bar of a door;}
}
{
}

\Verse{152}
{
    \zA{众顽童也有帮着打太平拳助乐的，也有胆小藏过一边的，也有立在桌上拍着手 乱笑、喝着声儿叫打的：登时鼎沸起来。}
}
{
    \eA{while Sao Hung and Ch'u Yo both laid hold of horsewhips, and they all rushed
        forward like a hive of bees. Chia Jui was driven to a state of exasperation;
        now he kept this one in check, and the next moment he reasoned with another,
        but who would listen to his words?}
}
{
}

\Verse{153}
{
    \zA{外边几个大仆人李贵等听见里边作反起来，忙都进来一齐喝住，问是何故，众 声不一，这一个如此说，那一个又如彼说。}
}
{
    \eA{They followed the bent of their inclinations and stirred up a serious
        disturbance. Of the whole company of wayward young fellows, some there were
        who gave sly blows for fun's sake; others there were who were not gifted
        with much pluck and hid themselves on one side; there were those too who
        stood on the tables, clapping their hands and laughing immoderately,}
}
{
}

\Verse{154}
{
    \zA{李贵且喝骂了茗烟等四个一顿，撵了出 去。}
}
{
    \eA{shouting out: "Go at it." The row was, at this stage, like water bubbling
        over in a cauldron, when several elderly servants,}
}
{
}

\Verse{155}
{
    \zA{秦钟的头早撞在金荣的板上，打去一层油皮，宝玉正拿褂襟子替他揉，见喝住 了众人，便命：“李贵，收书，拉马来!我去回太爷去!我们被人欺负了，不敢说别
        的，守礼来告诉瑞大爷，瑞大爷反派我们的不是，听着人家骂我们，还调唆人家打 我们。}
}
{
    \eA{like Li Kuei and others, who stood outside, heard the uproar commence
        inside, and one and all came in with all haste and united in their efforts
        to pacify them. Upon asking "What's the matter?" the whole bevy of voices
        shouted out different versions; this one giving this account,}
}
{
}

\Verse{156}
{
    \zA{茗烟见人欺负我，他岂有不为我的；他们反打伙儿打了茗烟，连秦钟的头也 打破了。}
}
{
    \eA{while another again another story. But Li Kuei temporised by rebuking Ming
        Yen and others, four in all, and packing them off. Ch'in Chung's head had,
        at an early period, come into contact with Chin Jung's pole and had had the
        skin grazed off.}
}
{
}

\Verse{157}
{
    \zA{还在这里念书么？”}
}
{
    \eA{Pao-yü was in the act of rubbing it for him, with the overlap of his coat,}
}
{
}

\Verse{158}
{
    \zA{李贵劝道：“哥儿不要性急，太爷既有事回家去了， 这会子为这点子事去聒噪他老人家，倒显的咱们没礼似的。}
}
{
    \eA{but realising that the whole lot of them had been hushed up, he forthwith
        bade Li Kuei collect his books. "Bring my horse round," he cried; "I'm going
        to tell Mr. Chia Tai-ju that we have been insulted. I won't venture to tell
        him anything else, but (tell him I will) that having come with all propriety
        and made our report to Mr. Chia Jui,}
}
{
}

\Verse{159}
{
    \zA{依我的主意，那里的事 情那里了结，何必惊动老人家。}
}
{
    \eA{Mr. Chia Jui instead (of helping us) threw the fault upon our shoulders.
        That while he heard people abuse us, he went so far as to instigate them to
        beat us;}
}
{
}

\Verse{160}
{
    \zA{这都是瑞大爷的不是，太爷不在家里，你老人家就 是这学里的头脑了，众人看你行事。}
}
{
    \eA{that Ming Yen seeing others insult us, did naturally take our part; but that
        they, instead (of desisting,) combined together and struck Ming Yen and even
        broke open Ch'in Chung's head. And that how is it possible for us to
        continue our studies in here?"}
}
{
}

\Verse{161}
{
    \zA{众人有了不是，该打的打，该罚的罚，如何等 闹到这步田地还不管呢？”}
}
{
    \eA{"My dear sir," replied Li Kuei coaxingly, "don't be so impatient! As Mr.
        Chia Tai-ju has had something to attend to and gone home, were you now, for
        a trifle like this, to go and disturb that aged gentleman, it will make us,
        indeed,}
}
{
}

\Verse{162}
{
    \zA{贾瑞道：“我吆喝着都不听。”}
}
{
    \eA{appear as if we had no sense of propriety: my idea is that wherever a thing
        takes place, there should it be settled;}
}
{
}

\Verse{163}
{
    \zA{李贵道：“不怕你老人 家恼我：素日你老人家到底有些不是，所以这些兄弟不听。}
}
{
    \eA{and what's the need of going and troubling an old man like him. This is all
        you, Mr. Chia Jui, who is to blame; for in the absence of Mr. Chia Tai-ju,
        you, sir, are the head in this school, and every one looks to you to take
        action.}
}
{
}

\Verse{164}
{
    \zA{就闹到太爷跟前去，连 你老人家也脱不了的。}
}
{
    \eA{Had all the pupils been at fault, those who deserved a beating should have
        been beaten, and those who merited punishment should have been punished!}
}
{
}

\Verse{165}
{
    \zA{还不快作主意撕掳开了罢！”}
}
{
    \eA{and why did you wait until things came to such a pass, and didn't even
        exercise any check?" "I blew them up," pleaded Chia Jui,}
}
{
}

\Verse{166}
{
    \zA{宝玉道：“撕掳什么？}
}
{
    \eA{"but not one of them would listen." "I'll speak out, whether you,}
}
{
}

\Verse{167}
{
    \zA{我必要 回去的！”}
}
{
    \eA{worthy sir, resent what I'm going to say or not,"}
}
{
}

\Verse{168}
{
    \zA{秦钟哭道：“有金荣在这里，我是要回去的了。”}
}
{
    \eA{ventured Li Kuei. "It's you, sir, who all along have after all had
        considerable blame attached to your name; that's why all these young men
        wouldn't hear you!}
}
{
}

\Verse{169}
{
    \zA{宝玉道：“这是为什 么？}
}
{
    \eA{Now if this affair is bruited, until it reaches Mr. Chia Tai-ju's ears,}
}
{
}

\Verse{170}
{
    \zA{难道别人家来得，咱们倒来不得的？}
}
{
    \eA{why even you, sir, will not be able to escape condemnation; and why don't
        you at once make up your mind to disentangle the ravelled mess and dispel
        all trouble and have done with it!" "Disentangle what?" inquired Pao-yü;}
}
{
}

\Verse{171}
{
    \zA{我必回明白众人，撵了金荣去！”}
}
{
    \eA{"I shall certainly go and make my report." "If Chin Jung stays here,"
        interposed Ch'in Chung sobbing,}
}
{
}

\Verse{172}
{
    \zA{又问李贵： “这金荣是那一房的亲戚？”}
}
{
    \eA{"I mean to go back home." "Why that?" asked Pao-yü. "Is it likely that
        others can safely come and that you and I can't?}
}
{
}

\Verse{173}
{
    \zA{李贵想一想，道：“也不用问了。}
}
{
    \eA{I feel it my bounden duty to tell every one everything at home so as to
        expel Chin Jung.}
}
{
}

\Verse{174}
{
    \zA{若说起那一房亲戚， 更伤了兄弟们的和气了。”}
}
{
    \eA{This Chin Jung," he went on to inquire as he turned towards Lei Kuei, "is
        the relative or friend of what branch of the family?" Li Kuei gave way to
        reflection and then said by way of reply:}
}
{
}

\Verse{175}
{
    \zA{茗烟在窗外道：“他是东府里璜大奶奶的侄儿，什么硬挣仗腰子的，也来吓我 们!璜大奶奶是他姑妈。}
}
{
    \eA{"There's no need whatever for you to raise this question; for were you to go
        and report the matter to the branch of the family to which he belongs, the
        harmony which should exist between cousins will be still more impaired."
        "He's the nephew of Mrs. Huang, of the Eastern mansion,"}
}
{
}

\Verse{176}
{
    \zA{你那姑妈只会打旋磨儿，给我们琏二奶奶跪着借当头，我 眼里就看不起他那样主子奶奶么。”}
}
{
    \eA{interposed Ming Yen from outside the window. "What a determined and
        self-confident fellow he must be to even come and bully us; Mrs. Huang is
        his paternal aunt! That mother of yours is only good for tossing about like
        a millstone, for kneeling before our lady Lien,}
}
{
}

\Verse{177}
{
    \zA{李贵忙喝道：“偏这小狗攮知道，有这些蛆嚼！”}
}
{
    \eA{and begging for something to pawn. I've no eye for such a specimen of
        mistress." "What!" speedily shouted Li Kuei, "does this son of a dog happen
        to know of the existence of all these gnawing maggots?"}
}
{
}

\Verse{178}
{
    \zA{宝玉冷笑道：“我只当是谁亲戚，原来是璜嫂子侄儿。}
}
{
    \eA{(these disparaging facts). Pao-yü gave a sardonic smile. "I was wondering
        whose relative he was," he remarked; "is he really sister-in-law Huang's
        nephew?}
}
{
}

\Verse{179}
{
    \zA{我就去向他问问。”}
}
{
    \eA{well, I'll go at once and speak to her." As he uttered these words,}
}
{
}

\Verse{180}
{
    \zA{说着便 要走，叫茗烟进来包书。}
}
{
    \eA{his purpose was to start there and then, and he called Ming Yen in, to come
        and pack up his books.}
}
{
}

\Verse{181}
{
    \zA{茗烟进来包书，又得意洋洋的道：“爷也不用自己去见他， 等我去找他，就说老太太有话问他呢。}
}
{
    \eA{Ming Yen walked in and put the books away. "Master," he went on to suggest,
        in an exultant manner, "there's no need for you to go yourself to see her;
        I'll go to her house and tell her that our old lady has something to ask of
        her. I can hire a carriage to bring her over,}
}
{
}

\Verse{182}
{
    \zA{雇上一辆车子拉进去，当着老太太问他，岂 不省事？”}
}
{
    \eA{and then, in the presence of her venerable ladyship, she can be spoken to;
        and won't this way save a lot of trouble?" "Do you want to die?" speedily
        shouted Li Kuei; "mind, when you go back,}
}
{
}

\Verse{183}
{
    \zA{李贵忙喝道：“你要死啊!仔细回去我好不好先捶了你，然后回老爷、 太太，就说宝哥儿全是你调唆。}
}
{
    \eA{whether right or wrong, I'll first give you a good bumping, and then go and
        report you to our master and mistress, and just tell them that it's you, and
        only you, who instigated Mr. Pao-yü! I've succeeded, after ever so much
        trouble, in coaxing them, and mending matters to a certain extent,}
}
{
}

\Verse{184}
{
    \zA{我这里好容易劝哄的好了一半，你又来生了新法儿! 你闹了学堂，不说变个法儿压息了才是，还往火里奔！”}
}
{
    \eA{and now you come again to continue a new plan. It's you who stirred up this
        row in the school-room; and not to speak of your finding, as would have been
        the proper course, some way of suppressing it, there you are instead still
        jumping into the fire."}
}
{
}

\Verse{185}
{
    \zA{茗烟听了，方不敢做声。}
}
{
    \eA{Ming Yen, at this juncture, could not muster the courage to utter a sound.}
}
{
}

\Verse{186}
{
    \zA{此时贾瑞也生恐闹不清，自己也不干净，只得委曲着来央告秦钟，又央告宝玉。}
}
{
    \eA{By this time Chia Jui had also apprehended that if the row came to be beyond
        clearing up, he himself would likewise not be clear of blame, so that
        circumstances compelled him to pocket his grievances and to come and entreat
        Ch'in Chung as well as to make apologies to Pao-yü.}
}
{
}

\Verse{187}
{
    \zA{先是他二人不肯，后来宝玉说：“不回去也罢了，只叫金荣赔不是便罢。”}
}
{
    \eA{These two young fellows would not at first listen to his advances, but
        Pao-yü at length explained that he would not go and report the occurrence,
        provided only Chin Jung admitted his being in the wrong. Chin Jung refused,
        at the outset,}
}
{
}

\Verse{188}
{
    \zA{金荣先 是不肯，后来经不得贾瑞也来逼他权赔个不是，李贵等只得好劝金荣，说：“原来 是你起的头儿，你不这样，怎么了局呢？”}
}
{
    \eA{to agree to this, but he ultimately could find no way out of it, as Chia Jui
        himself urged him to make some temporising apology. Li Kuei and the others
        felt compelled to tender Chin Jung some good advice: "It's you," they said,
        "who have given rise to the disturbance, and if you don't act in this
        manner,}
}
{
}

\Verse{189}
{
    \zA{金荣强不过，只得与秦钟作了个揖。}
}
{
    \eA{how will the matter ever be brought to an end?" so that Chin Jung found it
        difficult to persist in his obstinacy,}
}
{
}

\Verse{190}
{
    \zA{宝 玉还不依，定要磕头。}
}
{
    \eA{and was constrained to make a bow to Ch'in Chung. Pao-yü was, however, not
        yet satisfied,}
}
{
}

\Verse{191}
{
    \zA{贾瑞只要暂息此事，又悄悄的劝金荣说：“俗语说的：‘忍 得一时忿，终身无恼闷。’”}
}
{
    \eA{but would insist upon his knocking his head on the ground, and Chia Jui,
        whose sole aim was to temporarily smother the affair, quietly again urged
        Chin Jung, adding that the proverb has it: "That if you keep down the anger
        of a minute,}
}
{
}

\Verse{192}
{
    \zA{未知金荣从也不从，下回分解。}
}
{
    \eA{you will for a whole life-time feel no remorse." Whether Chin Jung complied
        or not to his advice is not known,}
}
{
}

\Verse{193}
{
    \zA{(本章完)}
}
{
    \eA{but the following chapter will explain.}
}
{
}
